\documentclass{article}
\usepackage[T1]{fontenc}
\usepackage{lmodern}
\usepackage{graphicx}
\usepackage{amsmath,amssymb}
\usepackage[a4paper,margin=2cm]{geometry}
\usepackage[absolute,overlay]{textpos}
\setlength{\TPHorizModule}{1mm}
\setlength{\TPVertModule}{1mm}
\usepackage[indonesian]{babel}
\usepackage{hyperref}
\setlength{\emergencystretch}{2em}
% unit: Halaman 7

\begin{document}

% === Halaman 7 (llm) ===

\begin{itemize}
    \item Contoh pada Caesar Cipher, penyerang hanya memiliki cipherteks berikut:
    
    KHOORZRUOG
    
    Penyerang tidak tahu kunci (pergeseran berapa huruf), hanya punya cipherteks itu saja.
\end{itemize}

Namun karena Caesar cipher hanya punya 25 kemungkinan pergeseran huruf, penyerang bisa mencoba semua pergeseran (\textit{brute force}).

Saat dicoba pergeseran 3 huruf ke kiri, diperoleh plainteks:

HELLOWORLD

Serangan berhasil menemukan kunci dan memecahkan cipherteks!

\end{document}
